\section*{Bab \ref{ch:b1:digestion-absorption} --- Pencernaan dan Penyerapan}

\begin{solution}{exo:b1:digestion-absorption:1}
Mulut: mengunyah, amilase liur. Kerongkongan: pengangkutan. Lambung:
asam, pepsin, simpanan dan campuran. Usus halus: empedu dan \omterm{def:b1:enzymes:enzyme}{enzim}
pankreas merampungkan \omterm{def:b1:digestion-absorption:nutrition}{pencernaannya}; \omterm{def:b1:digestion-absorption:nutrition}{penyerapan}. Usus besar: air dan
garam dipungut kembali, \omterm{def:b1:respiration-fermentation:fermentation}{fermentasi} oleh mikroba. Rektum: simpanan dan
pembuangan.
\end{solution}

\begin{solution}{exo:b1:digestion-absorption:2}
Pepsin dalam lambungnya pada pH 2; tripsin, kimotripsin dan
karboksipeptidase dari pankreasnya dalam usus halusnya pada pH 8;
peptidase tepi sikat pada pH 7. Sebagai zimogen, sehingga keduanya tak
mencerna \omterm{def:b1:cell-unit-of-life:cell}{sel} yang membuatnya; semuanya diaktifkan hanya dalam lumennya.
\end{solution}

\begin{solution}{exo:b1:digestion-absorption:3}
Glukosa: simpor natrium masuk, pembawa keluar, ke vena porta dan
hatinya. \omterm{def:b1:proteins:aminoacid}{Asam amino}: simpor natrium masuk, pembawa keluar, vena porta.
\omterm{def:b1:lipids:fattyacid}{Asam lemak}: difusi masuk dari sebuah misel, diesterkan kembali, dikemas
menjadi sebuah kilomikron, dieksositosis ke dalam lakteal dan limfanya.
\end{solution}

\begin{solution}{exo:b1:digestion-absorption:4}
Garam empedu mengemulsikan \omterm{def:b1:lipids:triglyceride}{lemak} menjadi tetesan semikrometer lalu
membawa hasil lipasenya dalam misel ke \omterm{def:b1:digestion-absorption:surface}{enterositnya}. Tanpa garam itu
\omterm{def:b1:lipids:triglyceride}{lemaknya} tinggal sebagai tetes besar yang hampir tak berpermukaan bagi
lipasenya, dan \omterm{def:b1:lipids:fattyacid}{asam lemak} yang terlepas tak dapat diantarkan ke
membrannya: \omterm{def:b1:lipids:triglyceride}{lemaknya} lewat tanpa tercerna.
\end{solution}

\begin{solution}{exo:b1:digestion-absorption:5}
$100/162 = \qty{0.62}{mol}$ satuan glukosa, sehingga sekitar
\qty{0.62}{mol} ikatan dan \qty{0.62}{mol} air (\qty{11}{g});
$620/180 = \qty{3.4}{mmol/min}$ glukosa yang terserap.
\end{solution}

\begin{solution}{exo:b1:digestion-absorption:6}
Asamnya melepaskan sekretin, yang membuat pankreasnya menggetahkan
bikarbonat yang menetralkan kimusnya; \omterm{def:b1:lipids:triglyceride}{lemak} dan peptidanya melepaskan
kolesistokinin, yang mengerutkan kantung empedunya (empedu masuk),
merangsang \omterm{def:b1:enzymes:enzyme}{enzim} pankreasnya, lalu menghambat pengosongan lambungnya,
sehingga duodenumnya menerima kimus hanya secepat ia sanggup
menetralkan lalu mencernanya.
\end{solution}

\begin{solution}{exo:b1:digestion-absorption:7}
Glukosanya dipekatkan delapan kali lawan landaiannya: \omterm{def:b1:membranes-transport:transporters}{pengangkutan
aktif}. Tanpa natrium pengangkutannya berhenti pada kesetimbangan:
serapannya bergandeng pada landaian natriumnya (\omterm{prop:b1:membranes-transport:secondary}{simporter}
natrium--glukosa), bukan langsung pada \omterm{def:b1:water-small-molecules:atp}{ATP}.
\end{solution}

\begin{solution}{exo:b1:digestion-absorption:8}
Kilomikron, selebar satu mikrometer, tak dapat melintasi dinding
kapilernya yang rapat tetapi dapat memasuki lakteal yang berujung
terbuka; limfanya bergabung dengan darahnya di lehernya, sehingga
\omterm{def:b1:lipids:triglyceride}{lemaknya} mencapai seluruh tubuhnya sebelum hatinya, yang bila tidak
akan menyerapnya seluruhnya lebih dahulu. Bila memasuki darah portanya
secara langsung, kilomikronnya akan menyumbat sinusoid hatinya lalu
menjenuhkan kesanggupannya mengolah \omterm{def:b1:lipids:triglyceride}{lemak} sesudah tiap hidangan.
\end{solution}

\begin{solution}{exo:b1:digestion-absorption:9}
Ia membasahi lalu mengapungkan seratnya bagi ruminasinya; bikarbonat
dan fosfatnya menyangga asam yang dihasilkan mikrobanya, sehingga
rumennya bertahan dekat pH 6,5; dan ia membawa urea dari darahnya,
yang diubah mikrobanya menjadi amonia lalu menjadi \omterm{def:b1:proteins:peptide}{proteinnya} sendiri,
sehingga mendaur ulang nitrogen yang bila tidak akan hilang dalam
kemihnya.
\end{solution}

\begin{solution}{exo:b1:digestion-absorption:10}
Tanpa amilase, protease, lipase dan bikarbonat pankreasnya: \omterm{def:b1:carbohydrates:polysaccharide}{patinya}
tercerna sebagian hanya oleh liur dan tepi sikatnya; \omterm{def:b1:proteins:peptide}{proteinnya}
dipotong pepsin tetapi tidak sampai ukuran yang dapat diserap;
\omterm{def:b1:lipids:triglyceride}{lemaknya} sama sekali tak tercerna lalu muncul dalam tinja yang
menggembung, pucat dan berminyak. Gula dan \omterm{def:b1:carbohydrates:glycosidic}{disakarida} (\omterm{def:b1:enzymes:enzyme}{enzim} tepi
sikatnya tetap ada), sebagian \omterm{def:b1:carbohydrates:polysaccharide}{pati}, serta peptida kecil masih dapat
ditangani; \omterm{def:b1:lipids:triglyceride}{lemak} dan sebagian besar \omterm{def:b1:proteins:peptide}{proteinnya} tidak.
\end{solution}

\begin{solution}{exo:b1:digestion-absorption:11}
Kotoran sekumnya membawa \omterm{def:b1:proteins:peptide}{protein} mikroba dan vitamin B yang dibuat
dalam sekumnya, yang terletak di luar usus halusnya sehingga tak dapat
dicerna pada lintasan pertamanya; ketika dimakan, keduanya melewati
lambung dan usus halusnya lalu terserap. Sapinya memfermentasi sebelum
lambungnya, sehingga mikrobanya dicerna di hilirnya tanpa lintasan
kedua mana pun. Sekum kudanya juga di hilirnya, dan kudanya tidak
memakan kotorannya maupun memiliki jalur lain: \omterm{def:b1:proteins:peptide}{protein} mikrobanya
hilang.
\end{solution}

\begin{solution}{exo:b1:digestion-absorption:12}
Mikrobanya memberi sapinya energi \omterm{def:b1:carbohydrates:polysaccharide}{selulosa} sebagai \omterm{def:b1:lipids:fattyacid}{asam lemak}, \omterm{def:b1:proteins:peptide}{protein}
yang dibuat dari rumput yang miskin dan bahkan dari urea, serta tiap
vitamin B; sapinya memberi mereka sebuah tong yang hangat, tersangga
dan tanpa udara yang dijaga penuh dan teraduk, serta seratus liter
liur sehari. Biaya rancangannya: sepersepuluh energinya pergi sebagai
metana, tak ada glukosa yang terserap sehingga hatinya harus membuat
semuanya dari propionat, \omterm{def:b1:digestion-absorption:nutrition}{pencernaannya} memakan tiga hari, dan tongnya
beserta isinya seberat seperlima hewannya --- seekor sapi tak dapat
berlari cepat, dan seekor singa dapat.
\end{solution}

\begin{solution}{pb:b1:digestion-absorption:1}
\textbf{1.} Serat \qty{5}{kg}, \omterm{def:b1:carbohydrates:polysaccharide}{pati} dan gula \qty{2}{kg}, \omterm{def:b1:proteins:peptide}{protein}
\qty{1.5}{kg}.
\textbf{2.} $0.7\times 5 + 2 = \qty{5.5}{kg}$.
\textbf{3.} $5500\times 17.5 = \qty{96}{MJ}$; ALA $0.75\times 96 =
\qty{72}{MJ}$.
\textbf{4.} $0.08\times 96 = \qty{7.7}{MJ}$; $7700/55 = \qty{140}{g}$
metana, sekitar \qty{200}{L}.
\textbf{5.} $0.15\times 5500 = \qty{825}{g}$.
\textbf{6.} Sebagian besar \omterm{def:b1:proteins:peptide}{protein} pakannya diuraikan mikroba dalam
rumennya lalu dibangun kembali sebagai \omterm{def:b1:proteins:peptide}{protein} mikroba; mikrobanya
mengalir terus ke abomasum dan usus halusnya, tempat mereka dicerna
seperti daging mana pun.
\textbf{7.} \qty{30}{\%} seratnya, \qty{1.5}{kg}: sebagiannya
difermentasi dalam kolonnya, sebagian besarnya dibuang.
\textbf{8.} \omterm{def:b1:proteins:peptide}{Protein} pakan yang lolos $0.4\times 1500 = \qty{600}{g}$,
tercerna \qty{360}{g}; mikroba \qty{825}{g}; total \qty{1185}{g}
$\times 17.5 = \qty{20.7}{MJ}$.
\textbf{9.} $72 + 20.7 = \qty{93}{MJ}$.
\textbf{10.} $72/93 = \qty{77}{\%}$ (sekitar \qty{60}{\%} begitu panas
\omterm{def:b1:respiration-fermentation:fermentation}{fermentasinya} dan energi sumbangan kolonnya diperhitungkan dalam
anggaran yang lebih penuh; menurut perhitungan ini, tiga perempat).
\textbf{11.} \qty{93}{MJ} kurang dari \qty{110}{}: sapinya harus makan
lebih banyak (\qty{12}{kg}), atau rumput yang lebih baik, atau menarik
dari \omterm{def:b1:lipids:triglyceride}{lemaknya}.
\textbf{12.} Propionat; \omterm{def:b1:biosyntheses-integration:gluconeogenesis}{glukoneogenesis} dalam hatinya.
\textbf{13.} $0.9\times 2000\times 17.5 = \qty{31.5}{MJ}$.
\textbf{14.} $0.7\times 1500\times 17.5 = \qty{18.4}{MJ}$.
\textbf{15.} $0.45\times 5 = \qty{2.25}{kg}$; ALA $0.75\times 2250\times
17.5 = \qty{29.5}{MJ}$.
\textbf{16.} $0.15\times 2250 = \qty{340}{g}$, hilang dalam tinjanya
(usus belakangnya berada di luar usus halusnya).
\textbf{17.} $31.5 + 18.4 + 29.5 = \qty{79}{MJ}$; bagian ALA $29.5/79 =
\qty{37}{\%}$.
\textbf{18.} Tidak: ia memerlukan $110/79\times 10 = \qty{14}{kg}$
rumput, dan mengaturnya dengan makan enam belas jam sehari lalu
melewatkan makanannya dua kali lebih cepat.
\textbf{19.} Sapi $93/10 = \qty{9.3}{MJ/kg}$; kuda $79/10 =
\qty{7.9}{MJ/kg}$.
\textbf{20.} Asupan kasar $10\,000\times 0.85\times 17.5 =
\qty{149}{MJ}$ (lignin dan abu dikecualikan); metana \qty{7.7}{MJ}
adalah \qty{5}{\%} darinya; seratus sapi: $100\times 140\times 365 =
\qty{5.1}{t}$ metana setahun.
\textbf{21.} Sapi $10\times 70/24 = \qty{29}{kg}$ bahan kering dalam
salurannya (ditambah sepuluh kali itu berupa air); kuda $10\times 36/24 =
\qty{15}{kg}$.
\textbf{22.} Pada asasnya seluruh \qty{340}{g}, senilai \qty{6}{MJ}
beserta \omterm{def:b1:proteins:aminoacid}{asam amino} esensialnya; isi sekum seekor kuda tidak terpisah
menjadi sebuah fraksi lunak yang kaya zat gizi sebagaimana isi sekum
seekor kelinci, dan seekor hewan setengah ton tak dapat memungut
kotorannya sendiri dari tanah dalam keadaan yang berguna.
\textbf{23.} Mikroba rumennya membuat \omterm{def:b1:proteins:peptide}{protein} dari nitrogen urea dan
karbon jerami, dan sapinya mencerna mikrobanya: ia hampir tak
memerlukan \omterm{def:b1:proteins:peptide}{protein} pakan. Mikroba kudanya berada di luar usus
halusnya, sehingga \omterm{def:b1:proteins:peptide}{proteinnya} hilang, dan kudanya harus menemukan \omterm{def:b1:proteins:aminoacid}{asam
aminonya} dalam makanannya sendiri.
\textbf{24.} Kudanya menyerap glukosa lalu menyimpan \omterm{def:b1:carbohydrates:polysaccharide}{glikogen} dalam
ototnya bagi lari cepat tanpa oksigen; sapinya tak menyerap satu pun,
membuat glukosa perlahan dari propionat, lalu memikul sebuah tong
seratus liter --- dibangun untuk merumput, bukan untuk berlari.
\textbf{25.} Sekitar tiga perempat energi termetabolis sapinya menurut
anggaran ini (\qtyrange{60}{70}{\%} dalam perhitungan yang penuh)
sebagai \omterm{def:b1:lipids:fattyacid}{asam lemak} atsiri, lawan sepertiga bagi kudanya;
\qty{9.3}{MJ} lawan \qty{7.9}{MJ} tiap kilogram rumputnya.
\end{solution}
